\documentclass[11pt]{amsart}
\usepackage[margin=1.1in]{geometry}
\usepackage{amsmath,amssymb,amsthm,booktabs,array,hyperref,xcolor,graphicx}
\hypersetup{colorlinks,linkcolor=blue!50!black,citecolor=blue!50!black,urlcolor=blue!50!black}
\newtheorem{theorem}{Theorem}[section]
\newtheorem{proposition}[theorem]{Proposition}
\newtheorem{lemma}[theorem]{Lemma}
\newtheorem{corollary}[theorem]{Corollary}
\theoremstyle{definition}
\newtheorem{definition}[theorem]{Definition}
\newtheorem{remark}[theorem]{Remark}
\newtheorem{question}[theorem]{Question}
\newcommand{\Z}{\mathbb{Z}}\newcommand{\Q}{\mathbb{Q}}\newcommand{\R}{\mathbb{R}}\newcommand{\A}{\mathbb{A}}\newcommand{\F}{\mathbb{F}}
\newcommand{\T}{\mathbb{T}}\newcommand{\Tr}{\operatorname{Tr}}\newcommand{\Aut}{\operatorname{Aut}}\newcommand{\Stab}{\operatorname{Stab}}
\newcommand{\Eis}{\mathrm{Eis}}\newcommand{\Gt}{G_2}\newcommand{\Ff}{F_4}
\title[The Hecke module of the compact $F_4$ at level 3]{The Hecke module of the compact $F_4$ at level $3$:\\ identification of its fifteen eigen-systems\\ and three new stable forms}
\author{Christopher Younge}\thanks{Working draft, not for circulation.}
\address{Philadelphia, PA}
\date{September 27, 2026}
\begin{document}
\begin{abstract}
The compact $F_4$ over $\Z$ has, at the parahoric level $Q_3$ attached to the first fundamental coweight, a class set of $15$ elements (Part~I). We compute the two Hecke operators $T(\omega_1^\vee)(2)$ and $T(\omega_4^\vee)(2)$ on it---the second by nine explicit isomorphisms and a forced completion---and match each of the eight joint eigen-systems, at the prime $2$, with an Arthur parameter built from forms on smaller groups: the lifts of $\Delta$ and of the newform in $S_{12}(\Gamma_0(3))$; the parameter $\textsf{3.8.a.a}[3]$ on $\mathrm{Sp}_6$, which carries the Eisenstein prime $41$; $\textsf{3.6.a.a}\oplus[4]$ on $\mathrm{Sp}_6$ paired with $\Delta$, which explains a congruence modulo $13$ as the Eisenstein congruence of \textsf{3.6.a.a}; two compact $G_2$-forms of weight $2\omega_2$ and level $3$, computed on $G_2$ and checked with both of its Hecke operators; and a compact $\mathrm{SO}(7)$-form of weight $(3,1,1)$ and level $3$ on the genus of $E_7$, found by a computation in which all three of its Hecke eigenvalues at $2$ agree with values predicted from the $F_4$ data, and whose Frobenius polynomial factors as $\textsf{3.6.a.a}$ times a stable genus-$2$ parameter of weight $\mathrm{Sym}^2\otimes\det^6$ with $\lambda(2)=-12$---the parameter predicted from $F_4$ before it was found. The two $G_2$-forms and the genus-$2$ parameter admit no endoscopic description of the expected shapes. Every congruence in the module is explained either as a transported classical Eisenstein congruence or as level-raising at $3$.
\end{abstract}
\maketitle
\section{Introduction}
In Part~I \cite{PartI} we computed the class set of the compact $F_4$ over $\Z$ at the parahoric level $Q_3$ (the stabiliser of a $6$-dimensional null subspace of the Albert lattice modulo $3$): fifteen classes, nine on the standard Albert lattice $J_\Z$ and six on the Elkies--Gross isotope $J_E$, and the Hecke operator $T(\omega_1^\vee)(2)$ on it, with integer spectrum
\[139230,\quad 8631^{(5)},\quad 15057^{(2)},\quad 3087^{(2)},\quad 2331,\quad 567^{(2)},\quad 207,\quad -585.\]
The Eisenstein primes $691$, $73$ and $41$ predicted by the mass formula were confirmed there at the form level, and the eigenvalues $8631$ and $15057$ were identified as the theta lifts of $\Delta$ and of the unique newform \textsf{3.12.a.a} in $S_{12}(\Gamma_0(3))$, which explains $691$ and $73$ as classical Eisenstein congruences transported to $F_4$. The present paper identifies the remaining five eigenvalues.

The method is elementary and, we think, of independent use. Gross's Satake dictionary expresses $T(\omega_1^\vee)(p)$ and $T(\omega_4^\vee)(p)$ as linear functions of the characters of the $26$- and $52$-dimensional representations of $F_4$ evaluated at the Satake parameter; for a parameter that factors through a subgroup $H\times H'\subset F_4$ these characters are polynomials in the Hecke eigenvalues of the pieces. Given a candidate shape, the two $F_4$-eigenvalues therefore either determine the eigenvalues of the pieces or overdetermine them. We enumerate shapes, solve, and then \emph{compute the predicted pieces on their own groups}: compact $G_2$ in weight $2\omega_2$, the non-split $\mathrm{SO}(5)$ on the genus of $A_5$, and the split $\mathrm{SO}(7)$ on the genus of $E_7$, with two or three Hecke operators each. A parameter is accepted only when every invariant available agrees. For $-585$ this produced a prediction---a genus-$2$ parameter with spinor eigenvalue $-12$ at $2$, paired with \textsf{3.6.a.a}---that was found, exactly, inside an $\mathrm{SO}(7)$-form some sessions later.

\begin{remark}[Convention]\label{conv}
Throughout this paper, to \emph{identify} an eigen-system with a parameter means that the Satake parameter at $2$, as seen by the two operators $T(\omega_1^\vee)(2)$ and $T(\omega_4^\vee)(2)$ (\S\ref{sec:T4}), agrees with that of the named parameter, and where a parameter is built from forms on a smaller group, that those forms have been computed there and their own Hecke eigenvalues at $2$ agree. No Hecke operator of $F_4$ at a prime other than $2$ has been computed, so every identification below is a match at one prime, with two invariants; we regard this as strong evidence and not as proof. At level one the identification of $\Theta(\Delta)$ has been verified at $5$ on $F_4$ itself by an exact computation of $T(\omega_1^\vee)(5)$ (\S\ref{sec:five}); on the $G_2$ side the identification of the $9/4$-form is verified at $5$ (\S\ref{sec:g2five}); the level-$Q_3$ predictions at $5$ of \S\ref{sec:g2five} remain to be tested.
\end{remark}

\subsection*{Acknowledgements} The computations were designed and carried out in collaboration with an AI system (Claude, Anthropic), which also drafted the text; all code and data are available (Appendix~\ref{app:data}).

\section{The class set and the first operator}\label{sec:setting}
\subsection{Level one (recollection)} The compact $\Ff$ over $\Z$ is $\Aut(J_\Z)$ for the integral Albert algebra $J_\Z=\mathrm{Her}_3(\mathcal O)$, and its level-one class set has two elements \cite{Gross-Z,EG}: $J_\Z$, with $|\Aut|=4180377600=6|W(E_8)|$, and the Elkies--Gross isotope $J_E$, with $|\Aut|=634023936=3|{}^3D_4(2)|$; the mass is $691/380414361600$. Their theta series are $E_{12}\pm(\text{const}/691)\Delta$, so the level-one Hecke module is $\mathbb{C}\mathbf 1\oplus\mathbb{C}\,\Theta(\Delta)$ \cite{Shan}, and the congruence mod $691$ is Ramanujan's. We enumerated the $2$-neighbours of both lattices for the two maximal parabolics (rank-one trace-zero points of $J/2J$, $69615$ of them, in five $\Aut$-orbits of sizes $48600,14175,6075,405,360$; and $6$-dimensional null subspaces for the Heisenberg parabolic) and classified each neighbour by explicit isomorphism. The Brandt matrices at $p=2$ are
\[T(\omega_4^\vee)(2)=\begin{pmatrix}1264320&7646400\\1159704&7751016\end{pmatrix},\qquad T(\omega_1^\vee)(2)=\begin{pmatrix}25830&113400\\17199&122031\end{pmatrix},\]
with eigenvalues $\{8910720,\,104616\}$ and $\{139230,\,8631\}$. Both cuspidal eigenvalues agree with the Satake-parameter prediction from $\tau(2)=-24$ for the lift of $\Delta$ via $\iota:\mathrm{Sp}_6\times\mathrm{SL}_2\to\Ff$ with the principal $\mathrm{SL}_2$ (for $T(\omega_1^\vee)$: $N_2-(1+2+\dots+2^5)(1+2^{11}-\tau(2))=139230-63\cdot1971=15057$ at level $3$ below, and $8631$ at level one). The $819$ rank-one elements of $J_E$ that define its relation to the Leech lattice \cite{EG} appear as the orbit of size $360$; each has exactly $72$ standard neighbours.

\subsection{Level $Q_3$} Let $Q_3$ be the parahoric at $3$ attached to $\omega_1^\vee$ (stabiliser of a $6$-dimensional null subspace of $J/3J$). The class set has $15$ elements: nine on $J_\Z$ with stabiliser orders $288$, $1296$, $1944$, $3024$, $15552$, $23328$, $72576$, $645120$, $1244160$, six on $J_E$ with $36$, $288$, $648$, $972$, $2304$, $9072$; and
\[\sum_x\frac1{|\Gamma_x|}=\frac{2068163}{52254720}=\frac{691}{380414361600}\cdot21789040,\]
$21789040$ being the number of $Q_3$-cosets. The matrix of $T(\omega_1^\vee)(2)$ (Appendix~\ref{app:matrix}) has row sums $139230$, is mass-weighted symmetric, and has integer spectrum
\[139230,\quad 8631^{(5)},\quad 15057^{(2)},\quad 3087^{(2)},\quad 2331,\quad 567^{(2)},\quad 207,\quad -585.\]
Here $E=41\cdot73\cdot691$, as the mass formula \cite[(1)]{PartI} predicts ($41\mid3^8-1$, $73\mid3^{12}-1$, $691$ from $\zeta(-11)$, with $5$ and $7$ cancelling), and the form-level test passes at exactly these primes: $691$ on the $\Delta$-lift eigenspace, $73$ on the $15057$-eigenspace, $41$ on the $2331$-eigenspace, and at no other prime (the eigenvalue coincidences at $2161,53,31,71,19,271,13,239$ all fail the test).

\section{The Satake dictionary and the classical lifts}
\subsection{Identification of the new eigenvalues}\label{sec:ident} The theta correspondence for $\Ff\times\mathrm{PGL}_2$ \cite{Shan} lifts a newform $f$ of weight $12$ and level $3$ to eigenvalue $N_2-63\,(1+2^{11}-a_2(f))$. The space $S_{12}(\Gamma_0(3))$ has a unique newform, \textsf{3.12.a.a}, with $a_2=78$, $a_3=-3^5$; it lifts to $15057$, and $1+p^{11}-a_p(f)\equiv0\pmod{73}$ for $p=2,5,7$: the prime $73$ on $\Ff$ is the classical Eisenstein congruence of \textsf{3.12.a.a}, transported, exactly as $691$ is Ramanujan's. The remaining eigenvalues are \emph{not} such lifts: $3087,2331,567$ would need $a_2=-112,-124,-152$, outside the Ramanujan bound $2\cdot2^{11/2}$, and $207,-585$ are not of the required shape at all.

The dictionary behind these formulas is $T(\omega_1^\vee)(p)=p^8\chi_{26}(s)-(1+p^4)$, where $s$ is the Satake parameter and $\chi_{26}$ the character of the $26$-dimensional representation; for a lift via $\mathrm{Sp}_6\times\mathrm{SL}_2$ one has $26=(14,1)\oplus(6,2)$ with $e_p$ principal in $\mathrm{Sp}_6$, which reproduces $139230$, $8631$ and $15057$ from $a_2=2049,-24,78$. The same dictionary through the subgroup $G_2\times\mathrm{SL}_2\subset F_4$, under which $26=(7,3)\oplus(1,5)$, with the $\mathrm{SL}_2$ principal, sends a $G_2$ eigenform with $T(\omega_1^\vee)(2)=t$ (so $\Tr V_7=(t+1)/8$) to
\[p^8\Big(\tfrac{t+1}{8}(p^{-1}+1+p)+(p^{-2}+p^{-1}+1+p+p^2)\Big)-(1+p^4)=112(t+1)+1967\quad(p=2).\]
The compact $G_2$ at level $3$ has trivial-weight cuspidal eigenforms with $t=9$ \cite{LP}, and $112\cdot10+1967=3087$; this is a coincidence: a trivial-weight $G_2$ form has Hodge weights $\{0,\pm1,\pm2,\pm3\}$ on the $7$, so $(7,3)\oplus(1,5)$ does not have the infinitesimal character of the trivial representation of $\Ff$, and the correct lifts through $G_2$ involve the weight $2\omega_2$ (\S\ref{sec:g2lift}). (The trivial-weight form is itself the $G_2$-lift of \textsf{3.6.a.a}: $\chi_7=t_{\textsf{3.6.a.a}}\chi_{[2]}+\chi_{[3]}=5/4$.) Two facts about the compact $G_2$ at $3$ are used below. First, at all three parahorics at $3$---the two maximal ones (stabilisers of a null line, resp.\ a null plane, modulo $3$) have class number $2$, stabiliser orders $\{432,36\}$ and $\{48,108\}$, mass $13/432$ and a single cuspidal eigen-system $(T_2,T_5)=(9,810)$; the Iwahori has class number $3$, mass $13/108$, and the same eigen-system with multiplicity $2$---no trivial-weight form with $T_2=-12$ exists, so $207$ is not a $G_2\times\mathrm{SL}_2$ lift with trivial $\mathrm{SL}_2$-parameter (which would need $96(t+1)+1263=207$, i.e.\ $t=-12$). Second, the dictionary predicts $T(\omega_1^\vee)(5)=5^8\big(\tfrac{811}{125}\cdot\tfrac{31}{5}+\tfrac{781}{25}\big)-626$ on a $G_2\times\mathrm{SL}_2$ lift of the $t=9$ form; the corresponding predictions for the parameters identified below are the test of Remark~\ref{conv}.

\subsection{Temperedness} With $\chi_{26}(s)=(\lambda+17)/256$, a unitary Satake parameter forces $|\chi_{26}(s)|\le26$, i.e.\ $\lambda\in[-6673,6639]$. The two identified $\mathrm{PGL}_2$-lifts violate this ($\chi_{26}=33.8$ and $58.9$), as non-tempered Arthur parameters must; the $G_2$-lift $3087$ ($\chi_{26}=12.1$) and the then-unidentified $2331,567,207,-585$ ($9.2,2.3,0.9,-2.2$) all satisfy it. The $41$-carrier is therefore consistent with being tempered, and if it is a lift from a smaller group the lifting parameter cannot involve a principal $\mathrm{SL}_2$ of the type that produced $8631$ and $15057$.
\section{The second Hecke operator}
\subsection{A second Hecke operator}\label{sec:T4} We computed $T(\omega_4^\vee)(2)$ on the $15$ classes. Its $8{,}910{,}720$ neighbours per class are indexed by the $69615$ rank-one points of $J/2J$ and $128$ lifts of each; for each of the five $\Aut(J_\Z)$-orbits of points, the stabiliser of a representative point permutes its lifts, and a count on the level-$3$ orbits is constant on a stabiliser-orbit of lifts, so one explicit isomorphism per stabiliser-orbit suffices: the $128$ lifts fall into $2,2,1,2,2$ orbits, i.e.\ \emph{nine} isomorphisms (not $640$), found and verified exactly as for $T(\omega_1^\vee)(2)$. This gives the nine standard rows; the $(E,Z)$ block follows by weighted symmetry, and the $(E,E)$ block is then \emph{forced}: the conditions $[T(\omega_1^\vee)(2),T(\omega_4^\vee)(2)]=0$, weighted symmetry and the row sums have a unique solution. Checks: all $225$ entries are non-negative integers; every row sums to $8{,}910{,}720$; every standard row splits $1{,}264{,}320$ standard $+\ 7{,}646{,}400$ exotic and every exotic row $1{,}159{,}704+7{,}751{,}016$, exactly the level-one Brandt matrix, although the exotic rows were not computed from neighbours. The joint eigen-systems are
\begin{center}\small
\begin{tabular}{@{}rrrl@{}}\toprule
$T(\omega_1^\vee)(2)$ & $T(\omega_4^\vee)(2)$ & mult. & identification\\\midrule
$139230$ & $8910720$ & $1$ & constant\\
$15057$ & $336870$ & $2$ & \textsf{3.12.a.a}$[6]\oplus[9]\oplus[5]$ (predicted $336870$)\\
$8631$ & $104616$ & $5$ & $\Delta[6]\oplus[9]\oplus[5]$ (predicted $104616$)\\
$3087$ & $23040$ & $2$ & $G_2\times\mathrm{SO}(3)$-lift, $G_2$ weight $2\omega_2$, $T_2=9$ (\S\ref{sec:g2lift})\\
$2331$ & $11916$ & $1$ & \textsf{3.8.a.a}$[3]$ via $\mathrm{Sp}_6\times\mathrm{SL}_2$ (predicted $11916$)\\
$567$ & $-14040$ & $2$ & $G_2\times\mathrm{SO}(3)$-lift, $G_2$ weight $2\omega_2$, $T_2=-27/2$ (\S\ref{sec:g2lift})\\
$207$ & $28800$ & $1$ & $\textsf{3.6.a.a}\oplus[4]$ on $\mathrm{Sp}_6$, $\Delta$ on $\mathrm{SL}_2$ (\S\ref{sec:gc})\\
$-585$ & $9000$ & $1$ & $\wedge^2_0F\oplus F[2]$, $F$ an $\mathrm{SO}(7)$-form of weight $(3,1,1)$ (\S\ref{sec:585})\\\bottomrule
\end{tabular}
\end{center}
With $T(\omega_4^\vee)(p)=p^{11}\chi_{52}(s_p)-p^8\chi_{26}(s_p)-p^6(p^4+1)$ and the restriction of the adjoint representation to $\mathrm{Sp}_6\times\mathrm{SL}_2$, $52=(21,1)\oplus(1,3)\oplus(14',2)$, the three identified $\mathrm{PGL}_2$-lifts were predicted before the computation and all three predictions hold exactly; in particular the identification of the $41$-carrier with \textsf{3.8.a.a}$[3]$ is now confirmed by two independent operators. Through $G_2\times\mathrm{SL}_2$, $52=(14,1)\oplus(1,3)\oplus(7,5)$, the eigenvalue $23040$ determines the adjoint character of the level-$3$ form of $G_2$ at $2$: $\chi_{14}(s_2)=7/64$, consistent with a tempered Satake parameter. All twelve congruences of the tables above hold for $T(\omega_4^\vee)(2)$ as well, including those modulo $5$ and $7$, which are therefore genuine.

\emph{Genus-$2$ shapes and the prediction for $-585$.} For the genus-$2$ shapes $\psi=\wedge^2_0F\oplus F\otimes g\oplus[1]\oplus F[2]\oplus g[2]$, the pair $(\chi_{26},\chi_{52})$ determines the normalised traces $(\operatorname{tr}F(s_2),\operatorname{tr}\wedge^2_0F(s_2))$ up to a quadratic ambiguity. For $567$ and $207$ no shape admits a unitary solution: \emph{neither is a genus-$2$ lift of this kind}. For $-585$ the shape with $g=\textsf{3.6.a.a}$ gives the unitary, exactly rational solution $\operatorname{tr}F(s_2)=-3\cdot2^{-7/2}$, $\operatorname{tr}\wedge^2_0F(s_2)=-11/16$, i.e.\ an \emph{integral} spinor eigenvalue $\lambda_F(2)=-6$ if $F$ has weight $\mathrm{Sym}^6\otimes\det^3$ (resp.\ $-12$ for $\mathrm{Sym}^2\otimes\det^6$); the other shapes give irrational or non-unitary values. We therefore predict: \emph{there is a genus-$2$ Siegel cusp form of level dividing $3$ and weight $\mathrm{Sym}^6\otimes\det^3$ (or $\mathrm{Sym}^2\otimes\det^6$) with spinor Hecke eigenvalue $-6$ (resp.\ $-12$) at $2$, and the $-585$-form is its lift with \textsf{3.6.a.a}.} We tested the level. Through $\mathrm{SO}(5)\sim\mathrm{PGSp}_4$, forms of paramodular level $3$ correspond to algebraic modular forms on the genus of the root lattice $A_5$ (even, quinary, half-discriminant $3$), which has a single class, with $\mathrm{SO}(A_5)\cong S_6$; the Siegel weights $\mathrm{Sym}^j\otimes\det^k$ correspond to the $\mathrm{SO}(5)$-representations of highest weight $((j+2k-6)/2,\,j/2)$, e.g.\ $(3,3)$ for $\mathrm{Sym}^6\otimes\det^3$. By the Weyl character formula the invariant spaces are $0$ for every candidate weight: \emph{there is no such form of paramodular level $3$}. At the smaller level $K'$ given by the stabiliser of an isotropic line in $A_5/3A_5$ (thirty non-radical isotropic lines in two orbits, with stabilisers of orders $36$ and $72$) the spaces are non-zero, and we computed the $2$-neighbour operator there. We realised $V_{(3,3)}$ as the $84$-dimensional Casimir eigenspace in $\mathrm{Sym}^3(\wedge^2V)$, found the fifteen $2$-neighbours of $A_5$ with explicit rational isometries onto $A_5$, and transported the level structure (checks: the representation is multiplicative on rational isometries, its image preserves the invariant subspaces, and on trivial weight $T_2$ has eigenvalues $15$ and $6$). In weight $V_{(3,3)}$ the space has dimension $3+1=4$ and
\[T_2\ \text{has eigenvalues}\quad 6,\quad -\tfrac34,\quad \tfrac32,\quad \tfrac32 .\]
Since the $p$-neighbour operator corresponds to the (minuscule) standard representation of the dual group $\mathrm{Sp}_4$, a tempered eigenform has $T_2=2^{3/2}\operatorname{tr}F(s_2)$, and the predicted $\operatorname{tr}F(s_2)=-3\cdot2^{-7/2}$ gives $T_2=-3/4$, which occurs. The prediction also fixes $\operatorname{tr}\wedge^2_0F(s_2)=-11/16$, and we tested it with the operator $T_P$ attached to totally isotropic planes of $A_5/2A_5$: there are $30$ such neighbours (all isometric to $A_5$; the enumeration also returns the $15$ ordinary neighbours, separated by their elementary divisors), $T_P$ commutes with $T_2$, and since its Satake transform lies in the span of $\chi_{\wedge^2_0}$ and $1$ with leading coefficient $2^{\langle\rho,\lambda\rangle}=4$, the constant function calibrates it: $T_P=4\chi_{\wedge^2_0}(s_2)-1$. On $V_{(3,3)}$ the joint eigen-systems are $(T_2,T_P)=(6,\tfrac{51}{8}),\ (-\tfrac34,\tfrac{21}{16}),\ (\tfrac32,-\tfrac{57}{8})^{(2)}$. \emph{The second test fails}: the form with $T_2=-3/4$ has $\operatorname{tr}\wedge^2_0F(s_2)=37/64$, not $-11/16$, so its match in the first invariant was a coincidence. Conversely, lifting each of the three genus-$2$ eigen-systems through the shape with $g=\textsf{3.6.a.a}$ gives $(T(\omega_1^\vee)(2),T(\omega_4^\vee)(2))=(711,5112)$, $(-261,252)$ and $(-585,16776)$: the last reproduces the first $\Ff$ eigenvalue exactly but not the second ($9000$), because the first $\Ff$ operator constrains the Satake parameter of $F$ only along a line. None of the genus-$2$ forms of weight $\mathrm{Sym}^6\otimes\det^3$ at level $K'$ is the one predicted; the predicted parameter is found in \S\S\ref{sec:585}--\ref{sec:stab}, in its $\mathrm{Sym}^2\otimes\det^6$ version, inside an endoscopic form on $\mathrm{SO}(7)$.

\section{Congruences in the module}\label{sec:cong}
\subsection{Congruences between cuspidal eigenforms} Two eigenspaces of $\T$ are congruent modulo $\ell$ when their saturated integer eigenlattices meet modulo $\ell$; this is a congruence of eigenvectors, hence of eigen-systems for the whole Hecke algebra. Applied to the $23$ rational eigenforms of the Niemeier genus, the test recovers Harder's congruence modulo $41$ between the eigenforms with $\lambda_2=21600$ and $-7920$ \cite{CL}, together with the propagation of the Bernoulli primes among the lifts and a dozen congruences modulo $11,13,17,19,23$ between pairs of eigenforms. Applied to the $\Ff$ module at level $Q_3$ (omitting $\ell=2,3$) it gives the three Eisenstein congruences of the theorem and the following cuspidal ones:
\begin{center}\small
\begin{tabular}{@{}p{4.2cm}lp{7.4cm}@{}}\toprule
eigenvalues & $\ell$ & remark\\\midrule
$15057,\ 8631$ & $17$ (two-dimensional) & the classical congruence $\Delta\equiv$\textsf{3.12.a.a}, $\tau(2)-a_2=-102=-6\cdot17$, transported\\
$8631,\ 3087$ & $11$ (two-dimensional) & the $3087$-forms are $G_2\times\mathrm{SO}(3)$-lifts of weight-$2\omega_2$ $G_2$-forms (\S\ref{sec:g2lift})\\
$8631,\ 207$ & $13$ & cyclotomic level-raising at $3$ ($13\mid3^3-1$); the $207$-form is not built from $\mathrm{PGL}_2$\\
$207,\ -585$ & $11$ & \\
$15057,3087$; $15057,567$; $8631,2331$; $3087,567$ & $5$ & confirmed by $T(\omega_4^\vee)(2)$ (\S\ref{sec:T4})\\
$2331,\ 567$ & $7$ & confirmed by $T(\omega_4^\vee)(2)$\\\bottomrule
\end{tabular}
\end{center}
The congruences modulo $11$ and $13$ between the lift of $\Delta$ and eigenforms that are not lifts are the $\Ff$ analogues of Harder's congruence: if the $3087$- and $207$-forms are lifts from $\mathrm{Sp}_6$ or $\mathrm{Spin}(9)$, these primes should divide the corresponding $L$-values.
\subsection{Level-raising at $3$} Every congruence between the $\Delta$-lift, which is old (it has a vector fixed by the hyperspecial subgroup), and a form new at $3$ should be a level-raising congruence: a prime $\ell$ at which two eigenvalues $e_i,e_j$ of the Frobenius at $3$ of $\psi_\Delta=\Delta[6]\oplus[9]\oplus[5]$ on the $26$ satisfy $e_i\equiv3e_j\pmod\ell$ without $e_i=3e_j$. Computing all such pairs in $\Q(\alpha_3(\Delta))$ gives the level-raising primes $5,7,11,13,17,19,41,83,137,181,691,\dots$ of $\psi_\Delta$ at $3$, and every cuspidal congruence of the $\Delta$-lift is of this kind:
\begin{itemize}
\item mod $17$ (with the \textsf{3.12.a.a}-lift): internal to the $\Delta$-block, Ribet's condition $\tau(3)^2\equiv(1+3)^2\,3^{10}$;
\item mod $11$ (with the $G_2$-lift): between the $\Delta$-block and a trivial block, $\tau(3)\equiv3^3+3^8\pmod{11}$;
\item mod $13$ (with the $207$-form) and mod $5$ (with the $2331$-form): \emph{cyclotomic}, $e_i/e_j$ a power of $3$ with $13\mid3^3-1$, resp.\ $5\mid3^4-1$, independent of $\tau(3)$; these are primes of the parahoric index $[K:Q_3]=21789040=2^4\cdot5\cdot7\cdot13\cdot41\cdot73$.
\end{itemize}
For the constant function the same cyclotomic condition holds at $5,7,13$, and the congruences do \emph{not} occur, because the level-one mass has $5,7,13$ in its denominator and they cancel in $E$ (\cite[Thm.~2.1]{PartI}); for the $\Delta$-lift the level-one congruence module contains only $691$, nothing absorbs $13$ or $5$, and the congruences occur. We therefore propose the following organising principle, which accounts for every congruence in the $\Ff$ level-$Q_3$ module that involves a known form: \emph{a congruence is either Eisenstein, at a prime of $E$, and then a transported classical Eisenstein congruence; or level-raising of an old form at $3$, at a prime of the parahoric index not absorbed at level one or at a prime where the old form's Satake parameter at $3$ satisfies a level-raising condition.}

\section{Searches through parameters of $\mathrm{Sp}_6\times\mathrm{SL}_2$ and $\mathrm{Spin}(9)$}
\subsection{An Arthur-parameter search}\label{sec:arthur} A parameter $\psi$ for $\Ff$ that is a sum of blocks $\pi[d]$ ($\pi$ cuspidal on $\mathrm{PGL}_n$, $[d]$ the $d$-dimensional representation of the Arthur $\mathrm{SL}_2$) contributes to the trivial-weight spectrum only if its infinitesimal character at $\infty$ is that of the trivial representation, which on the $26$-dimensional representation is the multiset $\{-8,\dots,8\}\cup\{-4,\dots,4\}$ (the principal $\mathrm{SL}_2$ acts on $26$ as $[17]\oplus[9]$). We enumerated every $\psi$ built from the trivial representation, the newforms of level $1$ and $3$ and weight $\le16$ (computed from $q$-expansions: of level $3$ they are \textsf{3.6.a.a}, \textsf{3.8.a.a}, two of weight $10$, \textsf{3.12.a.a}, three of weight $14$, two of weight $16$), and their symmetric squares, satisfying this constraint, and evaluated $T(\omega_1^\vee)(2)=256\chi_{26}(s_2)-17$ on each. There are $117$; all eigenvalues are $\equiv0\pmod9$, and $62$ of the $1479$ residues-$0$ slots in the tempered range are occupied, so a random integer has about a $4\%$ chance of being hit. The search reproduces $139230=[17]\oplus[9]$, $8631$ ($\Delta[6]\oplus[9]\oplus[5]$) and $15057$ (\textsf{3.12.a.a}$[6]\oplus[9]\oplus[5]$), and finds
\[2331:\qquad \psi=f[2]\oplus f[4]\oplus\mathrm{Sym}^2f\,[3]\oplus[5],\qquad f=\textsf{3.8.a.a}\ (a_2=6).\]
This $\psi$ is the image under $\mathrm{Sp}_6\times\mathrm{SL}_2\subset\Ff$ of the $\mathrm{Sp}_6$-parameter $f[3]$ with the Arthur $\mathrm{SL}_2$ embedded diagonally: $26=(14,1)\oplus(6,2)$ restricts to $\wedge^2_0(f[3])\oplus f[3]\otimes[2]=([5]\oplus\mathrm{Sym}^2f[3])\oplus(f[4]\oplus f[2])$. It is confirmed by the prime it carries: $1+2^7-a_2(f)=123=3\cdot41$, $41\mid3^8-1$, and $f\equiv1\oplus\chi^7\pmod{41}$ is the classical Eisenstein congruence of the weight-$8$ newform of level $3$; substituting $1\oplus\chi^7$ for $f$ in $\psi$ gives exactly the multiset of the principal parameter, so $\psi\equiv[17]\oplus[9]\pmod{41}$. \emph{The prime $41$ on the $2331$-form is the Eisenstein congruence of \textsf{3.8.a.a}, transported to $\Ff$}, as $691$ on $8631$ is Ramanujan's and $73$ on $15057$ is that of \textsf{3.12.a.a}. This also explains the theta-invisibility: the form is attached to weight $8$, not $12$.

The only parameter found for $567$ involves four distinct newforms ($\textsf{3.6.a.a}[4]\oplus\textsf{3.8.a.a}[2]\oplus\Delta[2]\oplus\textsf{3.16}[2]\oplus[5]\oplus[1]$); its image would contain four commuting $\mathrm{SL}_2$'s and a nontrivial commuting Arthur $\mathrm{SL}_2$, of total rank $5>4=\operatorname{rk}\Ff$, so it is not realisable and the hit is a coincidence. No parameter of this kind gives $207$ or $-585$. Hence $567^{(2)},207,-585$ do not come from forms on $\mathrm{PGL}_2$ of level dividing $3$: their parameters involve a stable form on $G_2$, a non-lifted Siegel form of genus $2$, or are stable on $\Ff$.

The cuspidal congruences are consistent with this. The congruence $\Delta\equiv\textsf{3.12.a.a}\pmod{17}$ is Ribet level-raising at $3$: $\tau(3)^2-(1+3)^2\,3^{10}=-2^7\cdot3^4\cdot5\cdot17$.

\subsection{Exhaustive genus-$2$ and triality sweeps} On the genus of $A_5$ the two levels at $3$ are the only ones (the quadratic space $A_5/3A_5$ has a one-dimensional radical and an anisotropic quotient of type $O^-(4,3)$: its $30$ isotropic lines lie over $10$ points), so we computed \emph{every} admissible genus-$2$ weight at both levels with both operators $T_2,T_P$ and lifted each eigen-system to $\Ff$ through its shape. Paramodular level: all spaces are $0$. Level $K'$: weights $(2,2),(3,2),(3,3),(4,1)$ of dimensions $3,2,4,3$ (weight $(1,1)$ is $0$). Two of the nine eigen-systems are Yoshida lifts, $(T_2,T_P)=(-\tfrac32,\tfrac34)=\mathrm{Y}(\textsf{3.8.a.a},\textsf{3.6.a.a})$ in weight $\mathrm{Sym}^4\otimes\det^3$ and $(-\tfrac{15}2,\tfrac{39}4)=\mathrm{Y}(\textsf{3.10.b},\textsf{3.6.a.a})$ in weight $\mathrm{Sym}^4\otimes\det^4$; as a check, their $\Ff$-lifts coincide, $(-585,3816)$, because they are built from the same three newforms. The other seven are not Yoshida or Saito--Kurokawa lifts of forms of level dividing $3$ and are presumably stable. \emph{No lift matches any of $(567,-14040)$, $(207,28800)$, $(-585,9000)$.}

We also searched the parameters through $\mathrm{Spin}(8)\supset\mathrm{SL}_2^4$, where $26=1+1+8_v+8_s+8_c$ and each $8$ is a sum of tensor products of pairs of the four factors: filling the factors with our newforms and the Arthur $\mathrm{SL}_2$ gives exactly two admissible parameters up to triality; the one built from \textsf{3.6.a.a}, \textsf{3.8.a.a}, \textsf{3.10.b} and $[2]$ again gives $(-585,3816)$, and the other nothing. The first operator thus returns $-585$ for several unrelated parameters---it sees only $\chi_{26}$, which is linear in the parameter's traces---while the second separates them, and none has $T(\omega_4^\vee)(2)=9000$. At this stage none of $567$, $207$, $-585$ was a lift from any family testable with forms of level dividing $3$ on the genera of $A_5$ and through $\mathrm{Spin}(8)$; they are resolved in \S\S\ref{sec:gc}, \ref{sec:g2lift} and \ref{sec:585}.

\subsection{A guess-and-check engine}\label{sec:gc} Any parameter through $\mathrm{Sp}_6\times\mathrm{SL}_2$ is determined by an $\mathrm{Sp}_6$-parameter $P$ and an $\mathrm{SL}_2$-parameter $Q$, and then $26=\wedge^2_0P\oplus P\otimes Q$ and $52=\mathrm{Sym}^2P\oplus\mathrm{Sym}^2Q\oplus\wedge^3_0P\otimes Q$; through $\mathrm{Spin}(9)\supset\mathrm{Spin}(5)\times\mathrm{Spin}(4)$ one has $26=1\oplus\wedge^2_0F\oplus g\otimes h\oplus F\otimes g\oplus F\otimes h$ and $52=\mathrm{Sym}^2F\oplus\wedge^2_0F\otimes g\otimes h\oplus\mathrm{Sym}^2g\oplus\mathrm{Sym}^2h\oplus F\otimes g\oplus F\otimes h$. Both $\Ff$-eigenvalues follow from the Satake eigenvalues of the pieces. We generated every $P$ of dimension $6$ built from newforms of level $1$ and $3$ (we also included those of level $9$---twists by the character of conductor $3$, CM forms by $\Q(\sqrt{-3})$, four quadratic pairs---although, the forms being fixed by a parahoric and hence Iwahori-spherical at $3$, they cannot occur; none enters any identification), Arthur blocks $[d]$, $f[d]$, $\mathrm{Sym}^2f[2]$, $\mathrm{Sym}^3f$, $f\otimes\mathrm{Sym}^2h$ and the genus-$2$ eigen-systems of \S\ref{sec:T4}, every $Q$, and also triple products through $\mathrm{Spin}(9)\supset\mathrm{SL}_2^3$, with the rank condition imposed. There are $351$ parameters with the infinitesimal character of the trivial representation. The engine reproduces $8631$, $15057$ and $2331$ exactly (the $G_2$-lift $3087$ lies outside it) and finds, uniquely,
\[207:\qquad\psi=\textsf{3.6.a.a}[4]\oplus[5]\oplus[1]\oplus(\textsf{3.6.a.a}\otimes\Delta)\oplus\Delta[4],\]
the image of the $\mathrm{Sp}_6$-parameter $\textsf{3.6.a.a}\oplus[4]$ with $\Delta$ on the $\mathrm{SL}_2$ factor; both $T(\omega_1^\vee)(2)=207$ and $T(\omega_4^\vee)(2)=28800$ are reproduced. The earlier searches missed it because it contains the four-dimensional tensor block $\textsf{3.6.a.a}\otimes\Delta$. Its congruence with the $\Delta$-lift modulo $13$ is now explained classically: $1+2^5-a_2(\textsf{3.6.a.a})=39=3\cdot13$ with $13\mid3^6-1$, so $\textsf{3.6.a.a}\equiv1\oplus\chi^5\pmod{13}$, and substituting this in $\psi$ turns $\textsf{3.6.a.a}[4]\oplus[5]\oplus[1]$ into $[9]\oplus[5]$ and $(\textsf{3.6.a.a}\otimes\Delta)\oplus\Delta[4]$ into $\Delta[6]$, the parameter of the $\Delta$-lift: \emph{$207\equiv8631\pmod{13}$ is the Eisenstein congruence of \textsf{3.6.a.a}, transported}, in the same way as $691$, $73$ and $41$. For $567$ and $-585$ no candidate reproduces both eigenvalues; twelve reproduce $T(\omega_1^\vee)(2)=-585$ alone, all containing \textsf{3.6.a.a}, which is again the weakness of the first operator.

\subsection{Exclusions for $567$ and $207$} Before their identification (\S\S\ref{sec:gc}, \ref{sec:g2lift}) we tested $567^{(2)}$ and $207$ against two further families of lifts; the exclusions remain of interest as uniqueness statements. (i) \emph{Stable $\mathrm{Sp}_6$-parameters} $F$ ($6$-dimensional, Hodge weights $\pm a/2,\pm b/2,\pm c/2$) through $\mathrm{Sp}_6\times\mathrm{SL}_2$: the shapes compatible with the infinitesimal character are $\wedge^2_0F\oplus F[2]$ with $(a,b,c)\in\{(11,5,3),(9,7,5)\}$ and $\wedge^2_0F\oplus F\otimes h$ with $h$ one of \textsf{3.6.a.a}, \textsf{3.8.a.a}, the two weight-$10$ forms, $\Delta$, \textsf{3.12.a.a}. Here $(\chi_{26},\chi_{52})$ leaves one free parameter; requiring the eigenvalues of $F$ on $\wedge^1,\wedge^2,\wedge^3$ to be rational integers after the arithmetic normalisation and the Satake parameter to be unitary leaves \emph{no} candidate for $567$ or $207$ (and a one-parameter family for $-585$, so this test does not decide $-585$). (ii) \emph{$\mathrm{SL}_3\times\mathrm{SL}_3$}: the only shapes compatible with the infinitesimal character are $\mathrm{ad}\,\pi\oplus\pi[3]\oplus\pi^\vee[3]$ and $\mathrm{ad}[3]\oplus[3]\otimes\pi\oplus[3]\otimes\pi^\vee$ with $\pi$ of Hodge weights $\{3,2,-5\}$, $\{5,-2,-3\}$, $\{4,3,-7\}$ or $\{7,-3,-4\}$; none of these is symmetric about its centre, so by Clozel's purity lemma no cuspidal $\pi$ on $\mathrm{PGL}_3$ has them, and the family is empty. Combined with \S\S\ref{sec:arthur}--\ref{sec:T4}: \emph{$567^{(2)}$ and $207$ are not lifts from $\mathrm{PGL}_2$ (level dividing $3$, weight $\le16$), $G_2\times\mathrm{SL}_2$, genus-$2$ Siegel forms, stable $\mathrm{Sp}_6$-forms with rational eigenvalues, or $\mathrm{PGL}_3$.} The parameter shapes not excluded are those through $\mathrm{Spin}(9)$ with a stable parameter of the group whose dual is $\mathrm{Spin}(9)$, lifts with non-rational Hecke fields, and stable parameters of $\Ff$. (Both were subsequently identified: $207$ in \S\ref{sec:gc} and $567$ in \S\ref{sec:g2lift}, through parameter shapes outside the families listed here.)

\subsection{Genus-$2$ candidates} If one of $567^{(2)},207,-585$ is a lift involving a genus-$2$ Siegel form $F$ (spinor Hodge weights $\pm a/2,\pm b/2$), the only shapes compatible with the infinitesimal character of $\Ff$ through $\mathrm{Sp}_6\times\mathrm{SL}_2$ or $\mathrm{Spin}(9)$ are $\psi=\wedge^2_0F\oplus F\otimes g\oplus[1]\oplus F[2]\oplus g[2]$ with
\begin{align*}(F,g)\in\{&(\mathrm{Sym}^2\!\otimes\det{}^3,\ \Delta),\ (\mathrm{Sym}^2\!\otimes\det{}^3,\ \textsf{3.12.a.a}),\ (\mathrm{Sym}^4\!\otimes\det{}^3,\ \text{wt }10),\\ &(\mathrm{Sym}^4\!\otimes\det{}^4,\ \textsf{3.8.a.a}),\ (\mathrm{Sym}^6\!\otimes\det{}^3,\ \textsf{3.6.a.a}),\ (\mathrm{Sym}^2\!\otimes\det{}^6,\ \textsf{3.6.a.a})\},\end{align*}
$F$ of level dividing $3$ (the shape $1\oplus\wedge^2_0F\oplus[3]\oplus[1]\oplus2F[2]$ has no solution). Of these six shapes exactly one occurs: the last, realised by the parameter $G$ of \S\ref{sec:stab} paired with \textsf{3.6.a.a}; the others are excluded by the two-operator test of \S\ref{sec:T4}.

\section{A second prime on $F_4$: the operator $T(\omega_1^\vee)(5)$ at level one}\label{sec:five}
The identifications of this paper are matches at the prime $2$ (Remark~\ref{conv}). At level one the theta-lift identification of the cusp form with $\Delta[6]\oplus[9]\oplus[5]$ predicts, from $\tau(5)=4830$ and the dictionary of \S\ref{sec:ident}, the eigenvalues
\[T(\omega_1^\vee)(5)\big|_{\Theta(\Delta)}=336{,}244{,}104,\qquad T(\omega_1^\vee)(5)\big|_{\mathbf 1}=N_5=191{,}040{,}038{,}280,\]
the second being $5$ times the number $38{,}208{,}007{,}656$ of $6$-dimensional sharp-null subspaces $S\subset J/5J$ (each has exactly five lifts, all Albert lattices). We have computed the operator exactly.

\emph{Method.} Direct enumeration of the $1.9\cdot10^{11}$ neighbours is out of reach; we compute the $\mathrm{Aut}(J_\Z)$-orbits of the subspaces $S$. The entry $a=\#\{J_\Z\text{-type neighbours of }J_\Z\}$ and the mass formula read
\[a=|\mathrm{Aut}(J_\Z)|\sum_{\mathrm{orbits}}\frac{n_Z(S)}{|\mathrm{Stab}(S)|},\qquad \sum_{\mathrm{orbits}}\frac1{|\mathrm{Stab}(S)|}=\frac{38208007656}{|\mathrm{Aut}(J_\Z)|}=\frac{75809539}{8294400},\]
with $n_Z(S)$ the number of lifts of $S$ isomorphic to $J_\Z$; the second identity certifies completeness: a list of pairwise distinct orbits whose masses sum to this number contains every orbit. Stabilisers were computed in two ways. For most $S$ the intersection lattice $M=J\cap J'$ (index $5^6$) has a small orthogonal group (orders $8$ to $18{,}432$, against $|O(J_\Z)|\approx9.7\cdot10^{28}$), and $\mathrm{Stab}(S)$ is the set of its elements that are integral Jordan automorphisms of $J_\Z$, verified exactly (the identity $(xA)^\#=x^\#A$ on all basis vectors and pairwise sums, and $\mathbf 1A=\mathbf 1$). When $O(M)$ is too large, we use the faithful action of $\mathrm{Aut}(J_\Z)$ on the $738$ vectors of norm $\le2$ (GAP confirms that our eleven generators generate a group of order exactly $4{,}180{,}377{,}600$): $\mathrm{Stab}(S)$ lies in the set stabiliser of the short vectors of $J$ contained in $M$, and is computed there by its action on subspaces of $\F_5^{27}$. Orbits are distinguished by isometry invariants of $M$ and, where these coincide, by exact isomorphism tests (four pairs of orbits have isometric intersection lattices and are separated only by the Jordan structure).

\emph{Orbits.} Uniform sampling ($10{,}238$ and $106{,}431$ subspaces) yields $41$ orbits with stabilisers of order $2$ to $192$, and $12$ orbits with large stabilisers ($96$ to $9600$); a final orbit, with stabiliser of order $110{,}592$ and expected frequency $\approx10^{-6}$ under uniform sampling, was found by sampling subspaces invariant under a fixed involution. The masses of these $54$ distinct orbits sum to exactly $75809539/8294400$, so the list is complete. All thirteen orbits with large stabilisers have $n_Z=5$. The result is
\[a=25{,}450{,}636{,}680,\qquad T(\omega_1^\vee)(5)=\begin{pmatrix}a&N_5-a\\[2pt] \tfrac{91}{600}(N_5-a)&N_5-\tfrac{91}{600}(N_5-a)\end{pmatrix},\]
with eigenvalues $N_5$ and
\[\lambda=336{,}244{,}104,\]
exactly the value predicted for $\Theta(\Delta)$; Ramanujan's congruence $\lambda\equiv N_5\pmod{691}$ holds. This confirms the identification of the level-one cusp form, and the Satake dictionary of \S\ref{sec:ident}, at a second prime on $F_4$ itself. The level-$Q_3$ predictions at $5$ (\S\ref{sec:g2five}) remain to be tested; the orbit computation above is the $J_\Z$-half of the class set of $F_4$ at the corresponding level-$5$ parahoric.

\section{The $G_2$-lifts of weight $2\omega_2$}
\subsection{The $G_2$-lifts of weight $2\omega_2$}\label{sec:g2lift} Through $G_2\times\mathrm{SO}(3)\subset\Ff$, with $26=(7,3)\oplus(1,5)$ and $52=(14,1)\oplus(1,3)\oplus(7,5)$, and the Arthur $\mathrm{SL}_2$ on the $\mathrm{SO}(3)$ factor, both $\Ff$-eigenvalues are \emph{linear} in the characters $(\chi_7,\chi_{14})$ of a $G_2$-parameter, so each open eigen-system determines them exactly. The infinitesimal character forces the $G_2$ Hodge weights $\{0,\pm3,\pm4,\pm7\}$, which by the Weyl dimension formula ($\dim=uv(u+v)(v-u)(2u+v)(u+2v)/120$ for Hodge type $(u,v)$) is the weight $2\omega_2$, of dimension $77$---not the trivial weight, and not $2\omega_1$ ($27$, type $(1,4)$). One finds: $3087\leftrightarrow(\chi_7,\chi_{14})=(5/4,7/64)$, $567\leftrightarrow(-25/16,329/128)$, both tempered, i.e.\ $G_2$-eigenvalues $T(\omega_1^\vee)(2)=8\chi_7-1=9$ and $-27/2$. We computed compact $G_2$ forms of weight $2\omega_2$ at all three levels at $3$ (the representation realised as the top Casimir component of $\mathrm{Sym}^2$ of the adjoint, built from the derivations of the octonions; the $126$ neighbour maps are octonion automorphisms; trivial weight reproduces $\{126,9\}$ and $\{126,9,9\}$; the operator preserves the invariant subspaces to $10^{-12}$). At the line-parahoric level the space has dimension $5$ and
\[T_2\ \text{has eigenvalues}\quad -\tfrac{27}2,\ -\tfrac{27}2,\ \tfrac94,\ 9,\ 9,\]
in exact correspondence, with multiplicities, with the five-dimensional part $3087^{(2)}\oplus2331\oplus567^{(2)}$ of the $\Ff$ module. The eigenvalue $9/4$ is the $G_2$-lift of \textsf{3.8.a.a} ($\chi_7=\chi_{[2]}t+t^2-1=13/32$, $\chi_{14}=0.8477$), whose $G_2\times\mathrm{SO}(3)$-lift reproduces \emph{both} $\Ff$-eigenvalues $(2331,11916)$, independently of its identification through $\mathrm{Sp}_6$ in \S\ref{sec:arthur}. We therefore identify $3087^{(2)}$ and $567^{(2)}$ as the $G_2\times\mathrm{SO}(3)$-lifts of the weight-$2\omega_2$ $G_2$-forms of level $3$ with $T_2=9$ and $-27/2$. The second invariants are checked with the second Hecke operator of $G_2$. Among the $10207$ lattices reached from the Coxeter order in two $\omega_1^\vee$-steps, classified by elementary divisors, those of relative position $(\tfrac14,\tfrac12,\tfrac12,1,1,2,2,4)$ number $2016=2^5\cdot63$ and form the $\omega_2^\vee$-double coset (the others: $1$, $126$, and $8064$ of type $2\omega_1^\vee$). Its Satake transform lies in the span of $\chi_{14},\chi_7,1$ with leading coefficient $2^{\langle\rho,\omega_2^\vee\rangle}=32$, and the constant function and the trivial-weight cusp form of level $3$ calibrate it: $T(\omega_2^\vee)(2)=32\,(\chi_{14}-\chi_7/4-1/2)$. On the five-dimensional weight-$2\omega_2$ space it commutes with $T_2$ and has eigenvalues $63/8$, $-45/2^{(2)}$, $315/4^{(2)}$ on the forms with $T_2=9/4,9,-27/2$, i.e.\ $\chi_{14}=217/256,\ 7/64,\ 329/128$. The first is the value predicted for the $G_2$-lift of \textsf{3.8.a.a}, and the other two are exactly the values required by the $\Ff$ eigen-systems of $3087$ and $567$. \emph{Both $\Ff$ eigenvalues of $3087$ and $567$ are thus reproduced by forms computed independently on $G_2$, with matching multiplicities.} The same method rules out a $G_2$ explanation of $-585$: its only admissible $G_2$ solution needs weight $2\omega_1$ with $T_2=99/4$ and the $\mathrm{SO}(3)$-parameter $\mathrm{Sym}^2$ of the CM form \textsf{9.4.a.a}; the weight-$2\omega_1$ spectra at the three levels are $\{0,81/2\}$, $\{0,81/2\}$ and $\{0^3,81/2^2,(-9\pm\sqrt{1161})/4\}$, and in any case a form fixed by the parahoric $Q_3$ is Iwahori-spherical at $3$, so every building block of its parameter has conductor dividing $3$ and level-$9$ forms cannot occur.

\subsection{The second prime on $G_2$}\label{sec:g2five} The operator $T(\omega_1^\vee)(5)$ on $G_2$ ($19530$ neighbour maps) was computed on the same weight-$2\omega_2$ space at the line-parahoric level; it commutes with $T_2$ (to $10^{-10}$) and the joint eigen-systems are
\[(T_2,T_5)=\Big(\tfrac94,\ \tfrac{14634}{25}\Big),\qquad\Big(9,\ -\tfrac{11502}{125}\Big)^{(2)},\qquad\Big(-\tfrac{27}{2},\ \tfrac{306}{5}\Big)^{(2)},\]
with $T(\omega_1^\vee)(p)=p^3\chi_7(s_p)-1$ in the unitary normalisation; the trivial-weight control at $5$ returns $\{19530,810\}$, the known values. The first eigen-system is a genuine second-prime test of the identification of the $9/4$-form as the $G_2$-lift of \textsf{3.8.a.a}: the lift predicts $\chi_7(s_p)=(p^{1/2}+p^{-1/2})t_p+t_p^2-1$ with $t_p=a_p/p^{7/2}$, and $a_5(\textsf{3.8.a.a})=390$ gives $T_5=14634/25$, exactly the computed value. The two stable forms have $\chi_7(s_5)=-11377/15625$ and $311/625$, both in the tempered range. Through $G_2\times\mathrm{SO}(3)$ these give the predictions
\begin{align*}T(\omega_1^\vee)(5)&=10439064\ \text{on the $3087$-eigenspace},\\ &=13407624\ \text{on the $567$-eigenspace},\\ &=23563224\ \text{on the $2331$-eigenspace},\end{align*}
the last agreeing exactly with the prediction of the $\mathrm{Sp}_6\times\mathrm{SL}_2$ route for \textsf{3.8.a.a}$[3]$ at $5$. These are the values an operator of $F_4$ at $5$ must return (Remark~\ref{conv}); its direct computation is out of reach, the number of $\omega_1^\vee$-neighbours at $5$ being $(5^{12}-1)(5^4+1)/4\approx3.8\cdot10^{10}$ per class.

\section{The eigenvalue $-585$ and the genus of $E_7$}
\subsection{Narrowing $-585$} Before the computation of \S\ref{sec:585}, three routes remained, and two were tested. (a) A genus-$2$ form $F$: over all pairings with newforms of level $1$ and $3$, the unique unitary solution with integral eigenvalues is $F$ of weight $\mathrm{Sym}^6\otimes\det^3$ with $\lambda_F(2)=-6$ and $\lambda_{\wedge^2}(2)=-352$ (or $\mathrm{Sym}^2\otimes\det^6$, $-12$, $-1408$), paired with \textsf{3.6.a.a}. It is absent from the genus of $A_5$ at both of its levels, so its local component at $3$ must be an Iwahori-spherical representation that does not transfer to the non-split $\mathrm{SO}(5)$; such a form is invisible to every compact inner form of $\mathrm{GSp}_4$ split at $2$, and deciding its existence requires holomorphic Siegel forms of level $\Gamma_0(3)$, which have not been tabulated (the computations of Bergstr\"om--Faber--van der Geer cover levels $1$ and $\Gamma_0(2)$). (b) A stable $\mathrm{Sp}_6$-parameter: the shapes compatible with $-585$ with integral, tempered eigenvalues require $\mathrm{SO}(7)$-weights $(1,1,0)$, $(2,1,0)$, $(2,2,0)$, $(3,0,0)$, $(2,2,2)$ or $(3,1,1)$. We computed compact $\mathrm{SO}(7)$-forms on the genus of $E_7$ (one class, $|\Aut|=2903040$; $63$ two-neighbours with explicit isometries; $364$ isotropic lines modulo $3$ in two orbits; trivial weight gives $T_2=63$ and $\{63,24\}$) in the first four weights, at level one and at the stabiliser of an isotropic line modulo $3$: all spaces are $0$ except one form of weight $(2,2,0)$ with $\lambda_1(2)=24$, where $-585$ would need $36$ or $39$. (c) A stable form on $\Ff$. The weights $(2,2,2)$, $(3,1,1)$ and the deeper levels of $\mathrm{SO}(7)$ at $3$ are computed in \S\ref{sec:585}, which resolves $-585$.

\subsection{$-585$ identified}\label{sec:585} The two largest weights, $(2,2,2)$ ($294$-dimensional, in $\mathrm{Sym}^2\wedge^3V$) and $(3,1,1)$ ($616$-dimensional, in $\mathrm{Sym}^2V\otimes\wedge^3V$), and the deeper levels at $3$---stabilisers of isotropic planes ($3640$, two orbits) and of isotropic $3$-spaces ($1120$, one orbit) modulo $3$---complete the $\mathrm{SO}(7)$ search. Weight $(3,1,1)$ has forms at both deeper levels, of dimensions $2$ and $1$, with $T_2$ giving $\lambda_1(2)=-60$: the first entry of the list of integral, tempered solutions for $-585$ in that shape, $(\lambda_1,\lambda_2,\lambda_3)=(-60,3264,-79872)$. We then built the other two Hecke operators at $2$ on the genus of $E_7$: the operator of the $630$ neighbours along totally isotropic planes modulo $2$, calibrated by the constant function as $T_P=16\chi_{\wedge^2_0}-5$, and the operator of the $1080$ neighbours along the $135$ maximal totally isotropic $3$-spaces, $T_3=2^{9/2}\chi_{\wedge^3_0}-2^{5/2}\chi_{\mathrm{std}}$; both commute with $T_2$. On the $(3,1,1)$-forms they give $T_P=9/2$, i.e.\ $\operatorname{tr}\wedge^2_0F(s_2)=19/32$, i.e.\ $\lambda_2=3264$; and $T_3=18$, i.e.\ $\lambda_3=-79872$. All three invariants of the Satake parameter at $2$ agree with the prediction, so
\[-585:\qquad\psi=\wedge^2_0F\oplus F[2],\]
with $F$ the $\mathrm{Sp}_6$-parameter of a compact $\mathrm{SO}(7)$-form of weight $(3,1,1)$ and level $3$,
with Hodge weights $\pm\tfrac{11}2,\pm\tfrac52,\pm\tfrac32$, reproduces both $\Ff$-eigenvalues $(-585,9000)$. Its Frobenius polynomial at $2$ shows that $F$ is endoscopic, $F=\textsf{3.6.a.a}\oplus G$ with $G$ a stable genus-$2$ form (\S\ref{sec:stab}). With this every one of the fifteen dimensions at level $Q_3$ is identified in the sense of Remark~\ref{conv}, each by two or more independent Hecke operators on a smaller group.

\section{Stability}
\label{sec:stab} From the three (resp.\ two) Hecke eigenvalues at $2$ one recovers the full Frobenius polynomial at $2$ of each new parameter, in the arithmetic normalisation, and its factorisation over $\Q$ decides endoscopy at the level of the Satake parameter: an endoscopic parameter's polynomial factors over the Hecke fields of its pieces, while a Galois group equal to the full Weyl group forbids every factorisation that the relevant endoscopic shapes would impose. As a control, the weight-$2\omega_2$ $G_2$-form with $T_2=9/4$ (the lift of \textsf{3.8.a.a}) gives $(X^2-96X+2^{15})(X^2-48X+2^{13})(X^2+220X+2^{14})$, with roots off the unitary circle, as a non-tempered lift must.
\begin{itemize}
\item \emph{The $G_2$-forms behind $3087$ and $567$.} Their Frobenius sextics $X^6-32X^5+18176X^4-720896X^3+\dots+2^{42}$ and $X^6+328X^5+58496X^4+7757824X^3+\dots+2^{42}$ are irreducible, with Galois group $D_6=W(G_2)$, and all roots of absolute value $2^7$. The Arthur shapes are excluded by temperedness; the $\mathrm{SL}_3$-shape $1\oplus\pi\oplus\pi^\vee$ would need $\pi$ of Hodge type $\{7,-3,-4\}$, excluded by Clozel's purity lemma; the $\mathrm{SO}(4)$-shapes need newforms of weight $4$ and $12$, and there are none of weight $4$ and level dividing $3$. \emph{No endoscopic parameter of these shapes reproduces them; we call the forms stable in this sense.}
\item \emph{The $\mathrm{SO}(7)$-form behind $-585$ is not stable.} Its Frobenius polynomial $X^6+60X^5+3264X^4+79872X^3+2^{11}\!\cdot3264X^2+2^{22}\!\cdot60X+2^{33}$ factors as $(X^2+48X+2^{11})(X^4+12X^3+640X^2+24576X+2^{22})$. The quadratic is \textsf{3.6.a.a} ($8a_2=-48$), and the quartic is a genus-$2$ parameter $G$ with $\lambda_G(2)=-12$ and $\operatorname{tr}\wedge^2_0G(s_2)=640/2^{11}-1=-11/16$: exactly the genus-$2$ form predicted for $-585$ in \S\ref{sec:T4}, in its weight-$\mathrm{Sym}^2\otimes\det^6$ version, paired with exactly the predicted newform. So the parameter of $-585$ is $\wedge^2_0G\oplus G\otimes\textsf{3.6.a.a}\oplus[1]\oplus G[2]\oplus\textsf{3.6.a.a}[2]$, the genus-$2$ shape of \S\ref{sec:T4}. The form $G$ is absent from the genus of $A_5$ because its local component at $3$ does not transfer to the non-split $\mathrm{SO}(5)$; it becomes visible as the $\mathrm{Sp}_4$-part of an endoscopic form on the split $\mathrm{SO}(7)$ of $E_7$.
\item \emph{The genus-$2$ parameter $G$ is stable in the same sense.} Its quartic is irreducible, with Galois group $D_4$, and tempered; a Yoshida lift would need newforms of weights $12$ and $4$, and a Saito--Kurokawa lift is non-tempered.
\end{itemize}
Thus the level-$Q_3$ module of $\Ff$ contains, besides lifts of classical newforms, three stable objects on smaller groups that had not been computed before: two compact $G_2$-forms of weight $2\omega_2$ and level $3$, and a stable parameter $G$ of $\mathrm{GSp}_4$ of weight $\mathrm{Sym}^2\otimes\det^6$ and level $3$ whose Hecke eigenvalue at $2$ was predicted from $\Ff$ before it was found. We have established $G$ only as a factor of the Frobenius polynomial of an $\mathrm{SO}(7)$-form; Arthur's classification for the split $\mathrm{SO}(7)$, in Ta\"ibi's extension to inner forms compact at infinity, should realise it as a holomorphic Siegel cusp form of level $3$, and we have not carried out that step.

\begin{question}
What are the Galois representations of the stable forms of \S\ref{sec:stab}, and which $L$-values do their congruences involve? What explains the remaining cuspidal congruences $567\equiv\{15057,3087\}\pmod5$, $567\equiv2331\pmod7$ and $207\equiv-585\pmod{11}$ in terms of these parameters?
\end{question}

\section{Concluding remarks}
Every one of the fifteen dimensions of the level-$Q_3$ module of $F_4$ is now matched with a parameter on a smaller group, in the sense of Remark~\ref{conv}, and every congruence in the module is accounted for by the organising principle of \S\ref{sec:cong}. Three of the parameters are new: two compact $G_2$-forms of weight $2\omega_2$ and level $3$, and a genus-$2$ parameter of weight $\mathrm{Sym}^2\otimes\det^6$ and level $3$ that is invisible on the non-split $\mathrm{SO}(5)$ and appears as the $\mathrm{Sp}_4$-part of an endoscopic form on the split $\mathrm{SO}(7)$. What remains is a confirmation at a second prime, which requires a Hecke operator of $F_4$ at $5$ (the neighbour count there is of the order of $4\cdot10^{10}$ per class, so a different method than direct enumeration is needed), the realisation of the genus-$2$ parameter as a holomorphic Siegel form through Arthur's classification, and the Galois representations of the three new forms.

\appendix
\section{The two Hecke matrices at level $Q_3$}\label{app:matrix}
Classes $1$--$9$ are on $J_\Z$ and $10$--$15$ on $J_E$, in the order of the stabiliser lists above. Entry $(i,j)$ counts the $\omega_1^\vee$-neighbours at $2$ of class $i$ lying in class $j$.
{\tiny\[\left(\begin{array}{rrrrrrrrrrrrrrr}
18270&3528&2268&1344&189&189&42&0&0&92232&10920&5040&3528&1344&336\\
15876&7011&0&2808&0&0&135&0&0&93744&9504&5400&3024&1296&432\\
15309&0&9279&0&567&675&0&0&0&91692&11340&3888&4212&2268&0\\
14112&6552&0&4851&0&0&252&63&0&84672&19656&7560&0&756&756\\
10206&0&4536&0&10521&378&0&0&189&81648&18144&0&9072&4536&0\\
15309&0&8100&0&567&1854&0&0&0&83592&13608&7776&8424&0&0\\
10584&7560&0&6048&0&0&1449&189&0&84672&12096&3024&0&4536&9072\\
0&0&0&13440&0&0&1680&8064&2646&0&94080&0&0&5880&13440\\
0&0&0&0&15120&0&0&5103&5607&0&90720&0&0&22680&0\\
11529&2604&1698&1008&189&129&42&0&0&99945&11214&5568&3582&1386&336\\
10920&2112&1680&1872&336&168&48&42&21&89712&21015&5280&3024&2376&624\\
11340&2700&1296&1620&0&216&27&0&0&100224&11880&8415&864&108&540\\
11907&2268&2106&0&567&351&0&0&0&96714&10206&1296&10413&3402&0\\
10752&2304&2688&576&672&0&144&21&42&88704&19008&384&8064&5679&192\\
10584&3024&0&2268&0&0&1134&189&0&84672&19656&7560&0&756&9387
\end{array}\right)\]}

The matrix of $T(\omega_4^\vee)(2)$, in the same ordering (row sums $8{,}910{,}720$):
\[\resizebox{\textwidth}{!}{$\left(\begin{array}{rrrrrrrrrrrrrrr}
851355&186480&122796&76800&13473&10233&2940&162&81&6196392&758136&343296&230760&94296&23520\\
839160&222876&81648&101304&6804&6804&5481&243&0&6223824&743040&359208&207792&84888&27648\\
828873&122472&220608&46656&26811&17199&1458&0&243&6164748&761724&311040&276372&120852&11664\\
806400&236376&72576&129348&6048&6048&5004&2142&378&6035904&926856&406728&145152&86940&44820\\
727542&81648&214488&31104&183618&17874&972&2187&4887&5820336&1013472&163296&428976&212544&7776\\
828873&122472&206388&46656&26811&31419&1458&0&243&6218856&795096&252720&280584&87480&11664\\
740880&306936&54432&120096&4536&4536&25830&5940&1134&5927040&1064448&311472&108864&72576&162000\\
362880&120960&0&456960&90720&0&52800&126324&53676&2903040&3743040&241920&0&336000&422400\\
349920&0&155520&155520&390960&12960&19440&103518&76482&2799360&3667680&0&311040&712800&155520\\
774549&172884&114162&71856&13473&9597&2940&162&81&6289722&762462&349584&230814&94914&23520\\
758136&165120&112848&88272&18768&9816&4224&1671&849&6099696&940392&342816&220752&114480&32880\\
772416&179604&103680&87156&6804&7020&2781&243&0&6292512&771336&401184&184896&70740&30348\\
778815&155844&138186&46656&26811&11691&1458&0&243&6231978&745038&277344&347454&137538&11664\\
754368&150912&143232&66240&31488&8640&2304&1200&1320&6074496&915840&251520&326016&161064&22080\\
740880&193536&54432&134460&4536&4536&20250&5940&1134&5927040&1035720&424872&108864&86940&167580
\end{array}\right)$}\]

\section{Data and code}\label{app:data}
All code (Python/NumPy, FLINT for exact characteristic polynomials), the class sets and Hecke matrices of $\Ff$ at levels $1$ and $Q_3$, the $G_2$ computations in weights $2\omega_1$ and $2\omega_2$, and the $\mathrm{SO}(5)$ and $\mathrm{SO}(7)$ computations on the genera of $A_5$ and $E_7$ are available from the author.

\begin{thebibliography}{99}
\bibitem{PartI} C.~Younge, Class sets and Eisenstein congruences on compact exceptional groups, working draft, 2026.
\bibitem{ATLAS} J.~H.~Conway, R.~T.~Curtis, S.~P.~Norton, R.~A.~Parker, R.~A.~Wilson, \emph{Atlas of Finite Groups}, Oxford, 1985.
\bibitem{BD} J.~Bergstr\"om, N.~Dummigan, Eisenstein congruences for split reductive groups, \emph{Selecta Math.} 22 (2016), 1073--1115.
\bibitem{BKK} T.~Berger, K.~Klosin, K.~Kramer, On higher congruences between automorphic forms, \emph{Math.\ Res.\ Lett.} 21 (2014), 71--82.
\bibitem{CDH} M.~Cheng, J.~Duncan, J.~Harvey, Umbral moonshine and the Niemeier lattices, \emph{Res.\ Math.\ Sci.} 1 (2014), Art.\ 3; arXiv:1307.5793.
\bibitem{CL} G.~Chenevier, J.~Lannes, \emph{Automorphic forms and even unimodular lattices}, Ergeb.\ Math.\ Grenzgeb.\ 69, Springer, 2019; neighbour tables at \url{http://gaetan.chenevier.perso.math.cnrs.fr/niemeier/niemeier.html}.
\bibitem{DS} P.~Deligne, J-P.~Serre, Formes modulaires de poids 1, \emph{Ann.\ Sci.\ ENS} 7 (1974), 507--530.
\bibitem{EG} N.~Elkies, B.~Gross, The exceptional cone and the Leech lattice, \emph{IMRN} 1996, 665--698.
\bibitem{GPR} S.~Garibaldi, H.~Petersson, M.~Racine, Albert algebras over $\Z$ and other rings, \emph{Forum Math.\ Sigma} 11 (2023).
\bibitem{Gross-Z} B.~Gross, Groups over $\Z$, \emph{Invent.\ Math.} 124 (1996), 263--279.
\bibitem{GN} B.~Gross, G.~Nebe, Globally maximal arithmetic groups, \emph{J.\ Algebra} 280 (2004), 375--398.
\bibitem{Gross-AMF} B.~Gross, Algebraic modular forms, \emph{Israel J.\ Math.} 113 (1999), 61--93.
\bibitem{Hohn17} G.~H\"ohn, On the genus of the moonshine module, arXiv:1708.05990.
\bibitem{HM} G.~H\"ohn, S.~M\"oller, Systematic orbifold constructions of Schellekens' vertex operator algebras from Niemeier lattices, \emph{J.\ London Math.\ Soc.}, to appear; arXiv:2010.00849.
\bibitem{LP} J.~Lansky, D.~Pollack, Hecke algebras and automorphic forms, \emph{Compositio Math.} 130 (2002), 21--48.
\bibitem{MS23} S.~M\"oller, N.~Scheithauer, Dimension formulae and generalised deep holes of the Leech lattice vertex operator algebra, \emph{Ann.\ of Math.} 197 (2023), 221--288.
\bibitem{MW} A.~Maloney, E.~Witten, Averaging over Narain moduli space, \emph{JHEP} 10 (2020) 187.
\bibitem{AJ} N.~Afkhami-Jeddi, J.~Cotler, T.~Hartman, A.~Maloney, A.~Tajdini, Free partition functions and an averaged holographic duality, \emph{JHEP} 01 (2021) 130.
\bibitem{Th76} J.~G.~Thompson, A conjugacy theorem for $E_8$, \emph{J.\ Algebra} 38 (1976), 525--530; P.~E.~Smith, A simple subgroup of $M?$ and $E_8(3)$, \emph{Bull.\ London Math.\ Soc.} 8 (1976), 161--165.
\bibitem{Witten07} E.~Witten, Three-dimensional gravity revisited, arXiv:0706.3359.
\bibitem{Mar17} K.~Martin, The Jacquet--Langlands correspondence, Eisenstein congruences, and integral $L$-values in weight 2, \emph{Math.\ Res.\ Lett.} 24 (2017), 1775--1795.
\bibitem{MW24} K.~Martin, S.~Wakatsuki, Mass formulas and Eisenstein congruences in higher rank, \emph{J.\ Number Theory} 257 (2024), 249--272.
\bibitem{Mazur} B.~Mazur, Modular curves and the Eisenstein ideal, \emph{Publ.\ Math.\ IHES} 47 (1977), 33--186.
\bibitem{NV} G.~Nebe, B.~Venkov, On Siegel modular forms of weight 12, \emph{J.\ reine angew.\ Math.} 531 (2001), 49--60.
\bibitem{Pollack-theta} A.~Pollack, Exceptional theta functions and arithmeticity of modular forms on $G_2$, preprint.
\bibitem{Shan} Y.~Shan, Level one algebraic modular forms on anisotropic $F_4$ and their applications, preprint.
\end{thebibliography}
\end{document}
